% Part: applied-modal-logics
% Chapter: epistemic-logic
% Section: truth-at-w

\documentclass[../../../include/open-logic-section]{subfiles}

\begin{document}

\olfileid[id]{aml}{el}{trw}

\olsection{Kebenaran pada Suatu Dunia}

Seperti pada logika modal normal, setiap model epistemik menentukan
!!{formula}s mana yang benar pada dunia mana di dalamnya. Kita menggunakan
notasi yang sama, yakni ``model~$\mModel{M}$ membuat !!{formula}~$!A$ benar
pada dunia~$w$,'' bagi gagasan dasar semantik relasional. Relasi ini
didefinisikan secara induktif dan identik dengan kasus modal normal untuk
semua operator nonmodal.

\begin{defn}\ollabel{defn:mmodels}
  \emph{Kebenaran !!a{formula}~$!A$ pada~$w$} dalam~$\mModel M$, dengan
  simbol $\mSat{M}{!A}[w]$, didefinisikan secara induktif sebagai berikut:
  \begin{enumerate}
  \tagitem{prvFalse}{\indcase{!A}{\lfalse}{Tidak pernah
    $\mSat{M}{\lfalse}[w]$}.}{}
  \tagitem{prvTrue}{\indcase{!A}{\ltrue}{Selalu
    $\mSat{M}{\ltrue}[w]$}.}{}
  \item $\mSat{M}{p}[w]$ jika dan hanya jika $w\in V(p)$
  \tagitem{prvNot}{\indcase{!A}{\lnot !B}{$\mSat{M}{\indfrm}[w]$ jika dan
    hanya jika $\mSat/{M}{!B}[w]$}.}{}
  \tagitem{prvAnd}{\indcase{!A}{(!B \land !C)}{$\mSat{M}{\indfrm}[w]$ jika
    dan hanya jika $\mSat{M}{!B}[w]$ dan $\mSat{M}{!C}[w]$}.}{}
  \tagitem{prvOr}{\indcase{!A}{(!B \lor !C)}{$\mSat{M}{\indfrm}[w]$ jika
    dan hanya jika $\mSat{M}{!B}[w]$ atau $\mSat{M}{!C}[w]$} (atau
    keduanya).}{}
  \tagitem{prvIf}{\indcase{!A}{(!B \lif !C)}{$\mSat{M}{\indfrm}[w]$ jika
    dan hanya jika $\mSat/{M}{!B}[w]$ atau $\mSat{M}{!C}[w]$}.}{}
  \tagitem{prvIff}{\indcase{!A}{(!B \liff !C)}{$\mSat{M}{\indfrm}[w]$ jika
    dan hanya jika $\mSat{M}{!B}[w]$ dan $\mSat{M}{!C}[w]$ keduanya berlaku,
    atau baik $\mSat{M}{!B}[w]$ maupun $\mSat{M}{!C}[w]$ tidak berlaku}.}{}
  \item\ollabel{defn:sub:mmodels-box}
    \indcase{!A}{\Knows_a !B}{$\mSat{M}{\indfrm}[w]$ jika dan hanya jika
    $\mSat{M}{!B}[w']$ untuk semua $w'\in W$ dengan $R_a ww'$}
  \end{enumerate}
\end{defn}

Di sinilah kita perlu memikirkan pembatasan terhadap relasi aksesibilitas.
Menurut klausa~\olref{defn:sub:mmodels-box}, !!a{formula}~$\Knows_a !B$ benar
pada~$w$ jika tidak ada $w'$ dengan $R_a ww'$. Ini merupakan klausa yang sama
dengan klausa dalam logika modal normal; jika suatu dunia tidak mempunyai
penerus, semua $\Box$-formula benar secara hampa di sana. Hal ini terasa
sangat berlawanan dengan intuisi jika $\Knows$ kita pahami sebagai
representasi \emph{pengetahuan}. Lagi pula, kita cenderung berpikir bahwa
\emph{tidak ada} keadaan tempat seorang agen dapat sekaligus mengetahui $!A$
dan $\lnot !A$.

Salah satu penyelesaiannya ialah memastikan bahwa relasi aksesibilitas kita
dalam logika epistemik selalu \emph{refleksif}. Ini secara kasar
berkorespondensi dengan gagasan bahwa dunia aktual konsisten dengan informasi
seorang agen. Logika epistemik pada umumnya menggunakan S5, tetapi logika lain
dapat menggunakan sistem yang lebih lemah, bergantung pada apa tepatnya yang
hendak direpresentasikan oleh relasi~$\Knows_a$.

\begin{figure}
  \begin{center}
    \begin{tikzpicture}[modal]
      \node[world] (w1) [label={[align=right]right:\mTrue{p}\\ \mFalse{q}}]{$w_1$};
      \node[world] (w2) [label={[align=right]right:\mTrue{p}\\ \mTrue{q}},
        above right=of w1]{$w_2$};
      \node[world] (w3) [label={[align=right]right:\mFalse{p}\\ \mFalse{q}},
        below right=of w1] {$w_3$};
      \draw[<->] (w1) to node {$a$} (w2);
      \draw[->,loop left] (w1) to node {$a,b$} (w1);
      \draw[<->] (w1) to [swap] node {$b$} (w3);
      \draw[->,loop above] (w2) to node {$a,b$} (w2) ;
      \draw[->,loop below] (w3) to node {$a,b$} (w3) ;
    \end{tikzpicture}
  \end{center}
  \caption{Suatu model epistemik sederhana.}
  \ollabel{fig:simple}
\end{figure}

\begin{prob}
  Tentukan pernyataan mana di antara pernyataan berikut yang berlaku dalam
  \olref[aml][el][trw]{fig:simple}:
  \begin{enumerate}
    \item $\mSat{M}{\lnot q}[w_1]$;
    \item $\mSat{M}{\Knows_a \lnot q}[w_1]$;
    \item $\mSat{M}{\Knows_b \lnot q}[w_1]$;
    \item $\mSat{M}{\Knows_b q \lor \Knows_b \lnot q}[w_2]$;
    \item $\mSat{M}{\Knows_a( \Knows_b q \lor \Knows_b \lnot q)}[w_2]$;
    \item $\mSat{M}{\EKnows_{\{a,b\}} \lnot q}[w_3]$;
    \end{enumerate}
\end{prob}

Setelah memberikan definisi dasar kebenaran pada suatu dunia, konsep semantik
lain dari logika modal normal, seperti validitas modal dan perikutan, langsung
terbawa serta dan diterapkan pada cara baru memahami penafsiran operator
modal ini.

Sekarang kita juga dapat memberikan syarat kebenaran bagi operator
pengetahuan bersama~$\CKnows_G$. Ingat dari \olref[sfr][rel][ops]{sec} bahwa
\emph{penutupan transitif}~$R^+$ dari suatu relasi~$R$ didefinisikan sebagai
\begin{align*}
  R^+ &= \bigcup_{0<n\in\Nat} R^n,
\intertext{dengan}
  R^1 &= R \text{ dan}\\
  R^{n+1} &= \Setabs{\tuple{x,z}}{\exists y(R^nxy\land Ryz)}.
\end{align*}
Kemudian, jika $G$ merupakan kelompok agen, kita mendefinisikan
$R_G=(\bigcup_{b\in G}R_b)^+$ sebagai penutupan transitif gabungan semua
relasi aksesibilitas agen.

\begin{defn}
Jika $G'\subseteq G$, kita menetapkan $\mSat{M}{\CKnows_{G'} !A}[w]$ jika dan
hanya jika, untuk setiap $w'$ sedemikian sehingga $R_{G'}ww'$,
$\mSat{M}{!A}[w']$.
\end{defn}

\end{document}
